% page 314 of the facsimile (image p-313 of the scan)
% block 154.9 x 246.3 mm ; left margin 8.0 mm
\pgc{-12.700mm}{0.000mm}{
\l{\cel{5}306}
\l{}
\l{\sou{kordono}. I. Kordeto qua esas un ek la elementi di kordo.}
\l{\cel{3}- II. Mikra treso ronda o plata, rubando plu o min larja,}
\l{\cel{3}uzata por rolo determinita. - III. (\sou{anat.)}\cel{2}Branchifuro}
\l{\cel{3}di la nervi. - IV. Lineo kontinua quan formacas kozo, se-}
\l{\cel{3}rio de kozi, alonge spaco determinita.}
\l{}
\l{\sou{koreo}. (\sou{patol.}) Morbo nervala qua produktas movi sukusoza.}
\l{\cel{3}- DEFS.}
\l{}
\l{\sou{koregono.}\cel{2}(\sou{zool.)}\cel{3}Genero di fisho qua vivas en aquo sen-}
\l{\cel{3}sala, familio di la salmonidei, qua vivas en la granda}
\l{\cel{3}lagi di la centro di Europa (Suisia, Germania, nordo-}
\l{\cel{3}regioni di Italia). -}
\l{}
\l{\sou{korekta}. Qua ne deviacas de la regulo, de la normo. DEFIRS.}
\l{}
\l{\sou{korelatar}\cel{2}(\sou{trans.}) Esar ligita ad altra termino tale ke un}
\l{\cel{3}ek la du advokas nemediate la altra. - DEFIRS.}
\l{}
\l{+\sou{korento}. (\sou{elektro).}\cel{2}Direciono determinita di la fluido}
\l{\cel{3}elektra, magnetala. (EFIS)}
\l{}
\l{\sou{korespondar}. (\sou{netrans., kun, pri}). I. Esar reciproke kon-}
\l{\cel{3}forma (l) pro interkonveno di sua naturo, di sua pro-}
\l{\cel{3}porcioni ; (2) pro interkonkordo di sentimenti. - II.}
\l{\cel{3}Esar reciproke komunikanta : (l) per paseyo komuna ;}
\l{\cel{3}(2) per kambio di letri.\cel{2}- DEFIRS.}
\l{}
\l{\sou{koriandro}. I. (\sou{bot.}) Planto aromata, qua formacas genero}
\l{\cel{3}ek la familio \textquotedbl{}ombeliferi\textquotedbl{}. - II. La semino sika di ta}
\l{\cel{3}planto, uzata en la koquo, en la farmacio-medikamenti.}
\l{\cel{3}- L.\cel{1}\sou{coriandrum}. - DEFIS.}
\l{}
\l{\sou{koribanto}. (\sou{epoki antiqua}). Sacerdoto di la diino Cibelo.}
\l{\cel{3}- DEFIRS.}
\l{}
\l{\sou{koridoro}. I. (\sou{olim)}\cel{2}Voyo tektoza cirkonde di fortifikajo.}
\l{\cel{3}II. (\sou{nia-epoke)}.\cel{2}Paseyo alonge plura chambri di un etajo,}
\l{\cel{3}di un apartmento, per qua onu povas forirar de oli. Pa-}
\l{\cel{3}seyo streta uzata kom forir-voyo, por irar de parto di}
\l{\cel{3}edifico ad altra parto. - DEFIRS.}
\l{}
\l{\sou{korifeo}. - Chefo di la koro, en la tragedio e la komedio}
\l{\cel{3}antiqua. - II. Chefo di koro kantala o dansala, en la}
\l{\cel{3}opero moderna. - III. (\sou{metaf}.) (\sou{kun nuanco di ironio)}}
\l{\cel{3}La persono qua maxime salias, unesma-ranga, en ulo.}
\l{\cel{3}- DEFIRS.}
\l{}
\l{\sou{korimbo.}\cel{2}(\sou{bot.}) Infloresenco di qua la pedunkli acesora,}
\l{\cel{3}departante de punti interdiferanta, elevas la flori a ni-}
\l{\cel{3}velo unika, tale ke formacesas quaza parasuno di qua la}
\l{\cel{3}radii esas neinteregala. - EFIS.}
\l{}
\l{\sou{korinto}. Bitbero sika, tre mikra, preske nigra, qua devenas}
\l{\cel{3}de Korinto, de la insuli Ioniana. - DFRS.}
}
